\documentclass[10pt,a4paper]{amsart}
\usepackage[margin=19mm]{geometry}
\usepackage{amsmath,amssymb,mathtools}
\usepackage[scr=rsfs,cal=euler]{mathalfa}
\usepackage{xfrac}
\usepackage[unicode,hidelinks,bookmarks=false]{hyperref}
\newcommand{\ie}{i.e.\ }
\newcommand{\cf}{cf.\ }
\newcommand{\resp}{respectively\ }
\newcommand{\Subsection}{\subsection}
\providecommand{\nobreakdash}{\nobreak\hbox{-}}
\setlength{\emergencystretch}{2em}
\multlinegap0pt
\usepackage{amssymb}
\usepackage{bm}
\usepackage{mathtools}
\usepackage{xcolor}
\newtheorem{proposition}{Proposition}
\newtheorem{lemma}{Lemma}
\newcommand{\A}{\mathscr{A}}
\newcommand{\Z}{\mathbb{Z}}
\renewcommand{\d}{%
  \relax\ifmmode\mathrm{d}\else\oldd\fi
}
\newcommand{\dR}{\mathrm{d}_{2,-1}}
\renewcommand{\dh}{%
  \relax\ifmmode\mathrm{d}_{1,0}\else\olddh\fi
}
\newcommand{\dv}{\mathrm{d}_{0,1}}
\title{Hamiltonian semisprays on Lie algebroids}
\author{Misael Avenda\~no Camacho, Jhonny Kama Mamani, Eduardo Velasco Barreras}
\date{Source: arXiv:2605.00483v1 (CC BY 4.0)}
\begin{document}
\maketitle

\noindent\emph{Source and licence.} Hamiltonian semisprays on Lie algebroids. Version arXiv:2605.00483v1 (CC BY 4.0), distributed by arXiv under the Creative Commons Attribution 4.0 International licence, http://creativecommons.org/licenses/by/4.0/, as recorded in the arXiv licence metadata for this exact version and shown on the abstract page https://arxiv.org/abs/2605.00483v1. Full licence notice: \texttt{licenses/arxiv-2605.00483-CC-BY-4.0.txt}.\par\smallskip

\noindent\textit{Editorial scope.} Counted unit: the article's lemma lemma:CohomD (source0.tex lines 754--762). The article's complete original proof of it is delivered with it (source0.tex lines 764--788, one \texttt{proof} environment of 25 lines). No mathematics is written, repaired or re-derived: the statement and the proof are the article's own lines, in the article's own notation, and no retained line is substituted or re-flowed.  Retained uncounted as the source-local context the proof needs: lines 689--695; lines 719--724.  Nothing was omitted from any retained span, so the package carries no omission record.  No retained line contains a citation, so the wrapper carries no bibliography and no bibliography entry.  No retained line contains a figure environment, an image-inclusion command or drawing code, so no image is delivered and the package declares no asset dependency.  Editorial formatting only, in addition: an AMS \texttt{amsart} preamble that restates the article's own declarations and macros used by the retained lines, namely the environments \texttt{proposition}, \texttt{lemma} on separate counters, together with the standard packages those lines need.  The retained environments print in this document as Proposition 1, Lemma 1 in order of appearance, whereas the article prints the same items with the numbering of its own full document.  The complete native source of the article as distributed in the arXiv e-print for this exact version is bundled unchanged as \texttt{sources/199566/source0.tex}.
\begin{align}
    \dv^2 &= 0, \label{eq:Cob02} \\
    \dh\dv + \dv\dh &= 0, \label{eq:Cob11} \\
    \dv\dR + \dR\dv + \dh^2 &= 0, \label{eq:Cob20} \\
    \dh\dR + \dR\dh &= 0, \label{eq:Cob3-1} \\
    \dR^2 &= 0 \label{eq:Cob4-2}
\end{align}

\begin{proposition}\label{vanishprop}
  Let $\A\to E$ be a Lie algebroid over a vector bundle $E\overset{\pi}{\to} M$,
  and $V\subset\A$ an involutive subbundle that is isomorphic as a Lie algebroid to the pullback bundle $p^*E$.
  Then, for any choice of an Ehresmann connection $\gamma$ for $V$ in $\A$, the cohomology of the cochain complex 
  $(\Gamma(\wedge^p V^\circ \otimes \wedge^{\bullet}H^\circ),\dv)$ vanishes in positive degrees, for every $p\in\Z$.
\end{proposition}

\begin{lemma}\label{lemma:CohomD}
  There exists a cochain complex $(\Gamma(\wedge^\bullet V^\circ|_M),D)$, where the coboundary operator $D$ is
  \[
    D\Theta_0 := (\dh\Theta)|_M,
  \]
  for $\Theta\in\Gamma(\wedge^\bullet V^\circ)$ the $\nabla$-parallel extension of $\Theta_0$.
  Furthermore, the $\nabla$-parallel extension gives a correspondence between 
  $D$-cocycles (resp. $D$-coboundaries) in $\Gamma(\wedge^\bullet V^\circ|_M)$ and $\d_\A$-closed (resp. $\d_\A$-exact) sections in $\Gamma(\wedge^\bullet V^\circ)$.
\end{lemma}

\begin{proof}
  Since $\Theta$ is horizontal and $\nabla$-parallel, we have $\dR\Theta=0$ and $\dv\Theta=0$. 
  By \eqref{eq:Cob20}, we have
  \[
    D^2\Theta_0 = (\dh^2\Theta)|_M = -(\dv\dR\Theta + \dR\dv\Theta)|_M = 0,
  \]
  proving that $D$ is a coboundary operator.
  Furthermore, our previous discussion implies that $\Theta$ is $\d_\A$-closed if and only if $\Theta$ is the $\nabla$-parallel extension of a $D$-cocycle $\Theta_0$.
  Now, if $\zeta_0\in\Gamma(\wedge^{p-1}V^\circ|_M)$ is such that $D\zeta_0 = \Theta_0$, then its $\nabla$-parallel extension $\zeta$ satisfies
  \[
    (\d_\A\zeta)|_M = (\dh\zeta + \dv\zeta)|_M = (\dh\zeta)|_M = D(\zeta_0) = \Theta_0 = \Theta|_M,
  \]
  implying that $\Theta=\d_\A\zeta$.
  Conversely, if $\Theta$ is $\d_\A$-exact, with $\Theta = \d_\A(\zeta+\xi)$ for $\zeta\in\Gamma(\wedge^{p-1}V^\circ)$ and $\xi\in\Gamma(\wedge^{p-2}V^\circ\otimes H^\circ)$,
  then by bidegrees, we must have $\dv\zeta=0$, $\dv\xi=0$, and $\dh\xi = 0$.
  The first condition imply that $\zeta$ is $\nabla$-parallel,
  and the second one, by Proposition \ref{vanishprop}, that $\xi=\dv\mu$ for $\mu\in\Gamma(\wedge^{p-2}V^\circ)$.
  By setting $\tilde{\zeta} := \zeta - \dh\mu$, we get from \eqref{eq:Cob11} and \eqref{eq:Cob20} that
  \begin{align*}
    \d_\A\tilde{\zeta} = \d_\A\zeta - \d_\A\dh\mu &= \d_\A\zeta - \dv\dh\mu - \dh^2\mu \\
    &= \d_\A\zeta + \dh\xi + \dv\dR\mu + \dR\dv\mu = \d_\A\zeta + \dR\xi = \d_\A(\zeta+\xi) = \Theta.
  \end{align*}
  Finally, since $\zeta$ is $\nabla$-parallel, and $\dh\xi=0$, it follows from \eqref{eq:Cob11} that $\tilde{\zeta}$ is also $\nabla$-parallel,
  showing that $\Theta$ is the $\nabla$-parallel extension of the $D$-coboundary with primitive $\tilde{\zeta}|_M$.
\end{proof}

\end{document}
